\chapter{Om duggan}
\label{ch:omduggan}

Duggan är skriftlig och genomförs med papper och penna. Tillåtna hjälpmedel är en fusklapp på ett
A4-blad, båda sidor, och en hexadecimal miniräknare. \textbf{Registerkartan delas ut med duggan} ---
indexen behöver inte memoreras, men att läsa kartan och veta vad den betyder behövs desto mer.

Byte-sekvenser, checksumberäkningar, registerpackning och SPI-transaktioner ska kunna göras
för hand.

\section*{Upplägg}

Duggan har två delar, ungefär lika stora:

\begin{booktable}{@{}lLr@{}}
\thead{Del} & \thead{Innehåll} & \thead{Poäng} \\ \midrule
A & Det egna protokollet: framedesign, parsning, adressering, felmodell och
  tillförlitlighetskedjan (\kapref{ch:frames}--\kapref{ch:bussar}). & 12 \\
B & CAN och drivrutinen: arbitrering, registerkartan, SPI-transaktionen, arkitekturen och
  begränsningarna (\kapref{ch:can}--\kapref{ch:avr}). & 13 \\
\end{booktable}

Godkänt kräver poäng i \emph{båda} delarna. Det går alltså inte att kompensera en tom del B med en
perfekt del A.

\section*{G- och VG-uppgifter}

\textbf{G-uppgifterna} kontrollerar grundläggande förståelse: bygga och läsa en frame, räkna en
checksumma, följa en parser genom en byte-ström, läsa registerkartan, packa en frame i register,
avkoda en SPI-transaktion, och förklara arbitrering.

\textbf{VG-uppgifterna} kräver resonemang: tillförlitlighetskedjan och dess konsekvenser,
analys av fel i en byte-ström, varför klibbiga statusbitar behövs, lagerarkitekturens gränser, och
vad konstruktionens begränsningar gör omöjligt för drivrutinen.

\section*{Övningsduggan}

Nästa kapitel är en fullständig övningsdugga i samma format och med samma poängfördelning som den
skarpa. Svaren står i kapitlet därefter --- med uträkningarna, inte bara resultaten.

Gör den med papper och penna, och läs svaren först efteråt.
